\documentclass[../report.tex]{subfiles}

\begin{document}

\subsection{Linux - Complete Pentest (Mesa)}

\subsubsection{Vulnerabilities}
Since the Mesa server uses a proprietary PHP server, no public exploits or CVE entries are known.
There are multiple vulnerabilities and avenues of compromising the system completely, which will be categorized here. 

\paragraph{CWE-548: Exposure of Information Through Directory Listing \autocite{mitre:cwe:548}}
The \textit{/backup} directory was indexed and accessible. This compromised usernames and password hashes. This could be prevented by preventing unintended directories from being indexed and restricting access.

\paragraph{CWE-759: Use of a One-Way Hash without a Salt \autocite{mitre:cwe:759}}
The compromised hashes were easily able to be cracked, in part because the hashes were calculated without employing a SALT. Including randomized data when hashing passwords would make it more difficult to crack passwords if the hashes were to be compromised \autocite{10331606}.

\paragraph{CWE-521: Weak Password Requirements \autocite{mitre:cwe:521}}
There seem to be no password requirements. Some passwords are regular english words, easily guessable or discoverable via Brute Force attack. More complex password guidelines would make the accounts less susceptible to breaches. Password guidelines can be implemented practically without hindering usability to an unreasonable extend, making user accounts more secure \autocite{10.1145/2891411}.

\paragraph{CWE-1391: Use of Weak Credentials \autocite{mitre:cwe:1391}}
Even if password requirements were introduced, the tendency of users to use passwords related to their personal life is dangerous. User passwords were found on publicly accessible webpages, and even if they were changed only slightly to adhere to complexity guidelines they could still be cracked via mutation of wordlists. Randomly generated passwords would be ideal, since they can't be inferred via OSINT. This would also solve the problem of reused passwords compromising multiple systems at once.

\paragraph{CWE-307: Improper Restriction of Excessive Authentication Attempts \autocite{mitre:cwe:307}}
 Online cracking targeting the login form on the Mesa homepage and SSH on the target machine was possible because multiple authentication attempts did not result in any restrictions. If repeated failed login attempts were detected, attacker IPs could be blacklisted.

\paragraph{CWE-308: Use of Single-factor Authentication \autocite{mitre:cwe:308}}
Since password discovery was possible, a complete compromise was possible via one simple avenue of attack. Requiring Multi-Factor-Authentication can improve security and prevent a password leak / crack from compromising accounts immediately.

\paragraph{CWE-89: Improper Neutralization of Special Elements used in an SQL Command \autocite{mitre:cwe:89}}
The user input intended for credentials gets not validated or escaped before being passed to an SQL query. Using Prepared Statements could be included in the PHP code running on the server to prevent users from manipulating SQL queries to gain illegitimate access \autocite{php:prep}.

\paragraph{CWE-23: Relative Path Traversal \autocite{mitre:cwe:23}}
URL parameters were not properly escaped, allowing for traversal into unintended directories, which allowed local file inclusion. This could be prevented by escaping special characters or validating the input, only allowing inputs matching a set of expected values to be parsed.

\paragraph{CWE-77: Improper Neutralization of Special Elements used in a Command \autocite{mitre:cwe:77})}
This concerns the same issue discussed in the previous vulnerability, the unvalidated user input in the Management-system's URL parameters. This allowed for the execution of shell commands on the target system, resulting in a reverse shell, compromising the entire system.

\paragraph{CWE-267: Privilege Defined With Unsafe Actions \autocite{mitre:cwe:267}}
A user was configured to be able to run a program with root privileges, but was able to run unsafe actions that were not intended, allowing for privilege escalation. This could be remedied by configuring the user rights correctly only for intended action, adhering to the Least Privilege principle \autocite{mitre:cwe:272}.

\newpage

\subsubsection{Risk assessment}

\paragraph{SQL Injection}
\begin{tabular}{l l}
\textbf{Vector}: & CVSS:4.0/AV:N/AC:L/AT:N/PR:N/UI:N/VC:H/VI:L/VA:L/SC:N/SI:N/SA:N\\
\textbf{CVSS v4.0 Score}: & 8.8 / High\\
\end{tabular}

The vulnerability works via Network (N) by accessing the webpage, has a Low Complexity (C), necessitates no Requirements (AT) or Privileges (P), or User Interaction (UI) on the server side.

This breaches Confidentiality (VC) initially, since usernames and password hashes are compromised upon login, which leads to compromise of the administrator accounts of the management system. The attacker could empty the database, therefore the Integrity and Availability is endangered, but the homepage itself remains unaffected.

\paragraph{Command Injection}
\begin{tabular}{l l}
\textbf{Vector}: & CVSS:4.0/AV:N/AC:L/AT:N/PR:N/UI:N/VC:H/VI:L/VA:L/SC:H/SI:N/SA:N\\
\textbf{CVSS v4.0 Score}: & 8.8 / High\\
\end{tabular}

This results in a reverse shell, which allows the attacker to read the everything the user (employee) has access to, possibly compromising sensitive data (e.g. user e-mails), which was explicitly forbidden from being examined in the penetration testing contract. Therefore the confidentiality breach has been assessed as High (H). Availability and Integrity have been classified as Low (L), since the webserver user \textit{www-data} accounts cannot alter source code or stop the server.

\paragraph{Offline Hash-Cracking / Online Brute Forcing (Webpage)}
\begin{tabular}{l l}
\textbf{Vector}: & CVSS:4.0/AV:N/AC:L/AT:N/PR:N/UI:N/VC:H/VI:N/VA:N/SC:H/SI:H/SA:H/CR:H/AU:Y\\
\textbf{CVSS v4.0 Score}: & 9.3 / Critical\\
\end{tabular}

Downloading the backup, and cracking the passwords (offline) or trying a wordlist with the found user account works via Network (N), has a Low Complexity (C), necessitates no Requirements (AT) or Privileges (P), or User Interaction (UI) on the server side.

This breaches Confidentiality (VC) initially, since usernames and password hashes are compromised. This additionally leads to a potential loss of Integrity and Confidentiality on the subsequent system (target machine), since the employee accounts use the same passwords.
These credentials allow login onto the system via SSH, which allows the attacker to read the everything the user (employee) has access to, possibly compromising sensitive data (e.g. user e-mails), which was explicitly forbidden from being examined in the penetration testing contract, therefore the Confidentiality breach has been categorized as High for subsequent systems. Availability and Integrity have been classified as Low (L), since the employee accounts can not alter source code or stop the server, but they could alter the user database.

\paragraph{Online Brute Forcing (Unix User)}
\begin{tabular}{l l}
\textbf{Vector}: & CVSS:4.0/AV:L/AC:L/AT:N/PR:N/UI:N/VC:H/VI:H/VA:H/SC:H/SI:H/SA:H\\
\textbf{CVSS v4.0 Score}: & 9.4 / Critical\\
\end{tabular}

Creating wordlists and running the Brute Force attack against the Unix system requires Local (L) connection via SSH, has a Low Complexity (C), necessitates no Requirements (AT) or Privileges (P), and no User Interaction (UI) on the server side.

Since this attack can compromise the root user \textit{zucc}, the exploit immediately and comprehensively breaches Confidentiality (VC), Integrity (VI), and Availability (VA), since root access can be obtained this way, endangering the totality of the host system, including subsequent systems like the webserver and database running on the system.

\paragraph{Privilege Escalation}
\begin{tabular}{l l}
\textbf{Vector}: & CVSS:4.0/AV:L/AC:L/AT:N/PR:L/UI:N/VC:H/VI:H/VA:H/SC:H/SI:H/SA:H\\
\textbf{CVSS v4.0 Score}: & 9.3 / Critical\\
\end{tabular}

The Privilege Escalation vulnerability can be exploited by logging in through SSH (Local (L)), has a Low Complexity (C), and necessitates no Requirements (AT). Low (L) Privileges (P) are necessary, since an account with the misconfiguration must be accessible to the threat actor. User Interaction (UI) on the server side is not required.
The damages are comparable to the prior exploit, since a root user is taken over with this attack as well.

\newpage
\paragraph{phpMyAdmin - Remote Code Execution}

\paragraph{Vulnerability - CVE-2018-12613}

\subparagraph{Working principle} \label{sec:5_pma_rce}
\enquote{The vulnerability comes from a portion of code where pages are redirected and loaded within phpMyAdmin, and an improper test for whitelisted pages} \autocite{mitre:cve:12613}.
The following explanation was adapted from \enquote{CHAMD5}'s explanation of the \enquote{phpMyAdmin 4.8.1 - (Authenticated) Local File Inclusion (1)} exploit \autocite{exploitdb:pma_lfi1}.

The code checking user input can be force to evaluate to TRUE, even if user input isn't included in the intended whitelist by prepending \textit{\%253f} to the malicious input. These double encoded URL characters lead to the check only considering the URL up to what the code expects to be the parameters, ignoring the File Inclusion string after the diversionary character (see \cref{lfi_check_pma}) \autocite{php:urldecode,php:mb-strpos}.
\begin{figure}[!htb]
\centering
\includegraphics[width=0.6\linewidth]{pentest/vulns/rce/check}
\caption{Code snippet validating user input via whitelist check}
\urlCaption{https://files.phpmyadmin.net/phpMyAdmin/4.8.1/phpMyAdmin-4.8.1-all-languages.zip}
\label{lfi_check_pma}
\end{figure}

The requested page gets validated and included with the malicious code still present (see \cref{lfi_code_pma}), abusing the improper loading of pages described in \textbf{CVE-2018-12613} \autocite{mitre:cve:12613}.

\begin{figure}[!htb]
\centering
\includegraphics[width=\linewidth]{pentest/vulns/rce/code}
\caption{Code snippet calling validation}
\label{lfi_code_pma}
\urlCaption{https://files.phpmyadmin.net/phpMyAdmin/4.8.1/phpMyAdmin-4.8.1-all-languages.zip}
\end{figure}

The script goes on to inject the command provided via command line parameter into the PHP session of phpMyAdmin, forcing it to load the code injected into the session via SELECT query by abusing the aforementioned file inclusion vulnerability (see \cref{lfi_exp_pma}).

\begin{figure}[!htb]
\centering
\includegraphics[width=\linewidth]{pentest/vulns/rce/exploit}
\caption{Code snippet executing exploit}
\label{lfi_exp_pma}
\urlCaption{https://www.exploit-db.com/exploits/44924}
\end{figure}

\subparagraph{Risk Assessment}
\begin{figure}[!htb]
\centering
\includegraphics[width=\linewidth]{pentest/vulns/rce/cvss.png}
\caption{CVSS v3.1 Score of CVE-2018-12613}
\urlCaption{https://nvd.nist.gov/vuln/detail/CVE-2018-12613}
\end{figure}

\begin{tabular}{l l}
\textbf{Vector}: & CVSS:3.1/AV:N/AC:L/PR:L/UI:N/S:U/C:H/I:H/A:H\\
\textbf{CVSS v3.1 Score}: & 8.8 / High\\
\end{tabular}

Any user with credentials can execute arbitrary code on the server. This vulnerability should urgently be patched by upgrading to a current version of \textit{phpMyAdmin}.  


\end{document}